\section{Simultaneous finite-horizon competitive factors}
\label{sec:finite-window-proof}

We prove Theorem~\ref{thm:finite-window} using the finite estimates proved
in the preceding two sections. No limit in $M$ or $r$ is taken. Throughout
this section $\ell=1+\log((M+2)/r)\ge1$, the certificate is $1/M$, and
$\Gamma$ denotes the finite-horizon quantities of
\eqref{eq:finite-factor}, not the sequence factors $C$.

\begin{proof}[Proof of Theorem~\ref{thm:finite-window}]
\noindent\textbf{Upper bound, including overlapping windows.}
Set $m=\min\{r-2,\lfloor M/2\rfloor\}$ and
$L=1+\log(M/(m+1))$. Since $r-1\ge r/2$ and
$\lfloor M/2\rfloor+1\ge M/2$,
\[
 \frac{M}{m+1}
 \le\max\left\{\frac{2M}{r},2\right\}
 \le\frac{2(M+2)}r.
\]
The last step uses $r\le M+2$ and covers full or overlapping windows.
Hence $L\le\ell+\log2\le2\ell$. Proposition~\ref{prop:finite-upper}
gives one counts-and-suffix map, independent of $k$, with all-start risk
at most $2^{19}Lg_M(k)\le2^{20}\ell g_M(k)$ for every $k$.
Stationary experiments are a subclass of arbitrary-start experiments,
so taking the infima over the same maps also gives
$\Gamma_{\rm stat}(M,r)\le\Gamma_{\rm all}(M,r)$.

\noindent\textbf{Lower bound and the small-logarithm case.}
Fix any allowable map $A$ and write
$G=\Gamma_{\rm stat}(A;M,r)\le\sqrt M$.
Lemma~\ref{lem:oracle-floor}, at $k=0$, implies $G\ge1/96$ even if
full history were available. If $\ell<64$, then
$1/96\ge2^{-30}\ell$, proving the desired bound. For $\ell\ge64$,
there is nothing to prove when $G\ge2^{-30}\ell$. Suppose instead
\begin{equation}\label{eq:fh-contradiction}
 G<2^{-30}\ell.
\end{equation}
As $r\ge2$ and $M\ge2$, $(M+2)/r\le M$. Thus
$\log M\ge\ell-1\ge63$, and $M>2^{80}$ because $80\log2<56<63$.
Choose the integer budget
\begin{equation}\label{eq:fh-budget}
 k=\left\lceil\max\{128r,M^{3/4}\}\right\rceil.
\end{equation}
This budget depends only on $(M,r)$; the same map $A$ is used at this
budget and at zero.

\noindent\textbf{Every finite lower-domain condition.}
Use the constants from Proposition~\ref{prop:finite-lower}, writing
$A_0=2^{12}$ for its numerical constant to distinguish it from the map:
\[
 c=2^{-16},\qquad t=ck/M,\qquad
 u_0=\lfloor ck/4\rfloor,\qquad
 L_-=\log\frac1{4A_0t}.
\]
The choice of $k$ gives $r\le k/128$, $k\ge M^{3/4}$, and
\[
 \frac M{k^2}\le M^{-1/2}<2^{-40}
   =\frac{A_0c^2}{2^{20}},\qquad
 ck\ge2^{-16}M^{3/4}>2^{44}\ge1024.
\]
Moreover
\[
 \frac rM=(1+2/M)e^{1-\ell}\le2e^{-63}<2^{-79}.
\]
Since $k\le128r+M^{3/4}+1$, this implies
\[
 \frac kM<2^{-72}+2^{-20}+2^{-80}<2^{-19}<\frac18.
\]
In particular $k$ is a legal budget, $k\le M/8$, and
$A_0t=k/(16M)<1/128<1/4$. Together with the original integer
$M\ge64,r\ge2$, these verify every condition in
\eqref{eq:finite-lower-domain}. The windows needed in that proof are
nonoverlapping since $r\le k/128\le M/1024$.
Also $ck\ge1024$ implies $u_0\ge ck/8$, so
\begin{equation}\label{eq:fh-prefactor}
 \rho:=\frac{u_0}{8Mt}\ge\frac1{64}.
\end{equation}

\noindent\textbf{The same map's clean risk.}
Let $S$ be the supremum of the clean risks on the stationary alternatives
in Proposition~\ref{prop:finite-lower}. All qualify at the same supplied
certificate, so $S\le G/\sqrt M$. For $x\ge0$, differentiation shows
that $(1+x)e^{-x/4}$ is maximized at $x=3$, with value
$4e^{-3/4}<4$. Thus $\ell\le1+\log M\le4M^{1/4}$ and
\eqref{eq:fh-contradiction} gives the finite inequalities
\begin{equation}\label{eq:fh-clean}
 S<2^{-28}M^{-1/4}\le2^{-21}M^{-1/4}\le t/32.
\end{equation}
The last step uses $k\ge M^{3/4}$ and $c/32=2^{-21}$.
No asymptotic small-order estimate or distinct clean estimator is used.

\noindent\textbf{A logarithm uniform in both parameters.}
Since $4A_0c=1/4$, $L_-=\log(4M/k)$. Put
$x=\max\{128r,M^{3/4}\}\ge1$; then $k=\lceil x\rceil\le2x$, and
\[
 L_-\ge\min\left\{\log\frac M{64r},\ \log(2M^{1/4})\right\}.
\]
The first term equals
$\ell-1-\log(1+2/M)-\log64$ and is at least
$\ell-2-\log64\ge\ell-8\ge\ell/8$, using
$\log(1+2/M)\le1$, $\log64<6$, and $\ell\ge64$.
The second term is at least $(\ell-1)/4\ge\ell/8$.
Consequently $L_-\ge\ell/8$.

\noindent\textbf{Common-kernel comparison and infimum.}
Proposition~\ref{prop:finite-lower} supplies a qualifying stationary null
and one strict source-independent kernel, preserving the actual final pair,
with
\[
 R_{\rm null}\ge\rho\left[\left(t/16-S\right)L_-\right]_+.
\]
Its full proof above includes source legality, mixture normalization,
window matching, and pointwise budget checks. Here the bracket is positive
by \eqref{eq:fh-clean}. Since $k\ge M^{3/4}\ge\sqrt M$, one has
$g_M(k)=k/M$. Therefore
\[
 G\ge\frac{R_{\rm null}}{g_M(k)}
 \ge\frac{(1/64)(t/32)(\ell/8)}{k/M}
 =2^{-30}\ell,
\]
contradicting \eqref{eq:fh-contradiction}. The lower bound thus holds for
every map and passes to its infimum without requiring attainment.
The witnesses are stationary and hence are also available in the all-start
class. This completes the finite-horizon proof.
\end{proof}

For clarity about the sequence consequences, if $r_M/M$ has positive
liminf, then $\ell_M$ is bounded eventually, and finitely many earlier
horizons have finite factor; the explicit suffix maps give a single
uniform family. If that liminf is zero, $\ell_M$ diverges on a subsequence,
so the finite lower forbids a uniform factor for any family. For fixed
$\alpha$, the finite bounds give the logarithmic order, while the original
fixed-$\alpha$ proof in Proposition~\ref{prop:lower-consequences} retains
its stronger $2^{-26}\log(1/(32\alpha))$ lower constant. Neither argument
claims a logarithmic risk lower at every individual budget.
